\chapter{29}

\section[«Sunnewaid», Bettmeralp (CH)]{\emph{«Sunnewaid»}, Bettmeralp (CH)} 

Es war frischer geworden. Eine vage Andeutung von Herbst lag in der
Luft. Louis blieb auf dem schmalen Pfad stehen. Er blickte weit über das
Tal und stellte sich breitbeinig hin.
Das Handy hatte ihm den Weg zuverlässig gewiesen. Hier musste es
irgendwo sein. Er schaute sich um. Tatsächlich. Prompt ertönte die
weibliche Stimme aus dem Gerät.

«Sie haben Ihr Ziel erreicht.».

Louis atmete auf. 

In überschaubarer Weite hörte er gedämpft Musik. Klavierklänge legten einen
feinen Klangboden über die Umgebung. Vor einem grosszügigen Holzhaus mit
angebauter Scheune sah er eine Frau auf einer Bank liegen. Ihr Kopf war
auf eine Erhöhung gebettet. Den rechten Arm liess sie neben der Bank über
dem Boden baumeln. Ihre Augen schienen geschlossen. Aus der Distanz
konnte es Louis nicht mit Sicherheit sehen. Die nächsten Häuser fanden
sich in allen Himmelsrichtungen und weiten Abständen.
Einige Schritte näher hörte Louis eine männliche Stimme singen. Der
Gesang hatte etwas Gelassenes und passte zum Musikstil. Er glaubte,
französische Worte zu hören. Das Lied startete zum zweiten Mal. Jetzt
verstand er den Anfang des Refrains. \emph{«Je te laisserai des mots»}\footnote{«Ich werde dir Worte da lassen.»}.

\qrblock{Patrick Watson \emph{«Je te laisserai des mots»}}{media/QR-Codes/026_Je_te_laisserai_des_mots_Patrick_Watson_Spotify.png}

Er blieb stehen. Knapp versteckt hinter dem Birnbaum, dem einzigen Baum
um das Haus, verfolgte er die Szene.
Von hier aus hatte er einen guten Überblick.
Kaum geendet, startete der Song von Neuem. Beim dritten Mal kicherte
Louis leise. Endlosschlaufe. Kannte er. Wenn es zu schön war, um
aufzuhören. Wenn am Ende die Leere lauerte.

Wenige Meter abseits auf einem Schotterweg, der zum Haus führte, ruhte
ein in die Jahre gekommener, silbergrauer BMW Turbo mit schwarzem
Streifen. Der zog sich geradlinig am unteren Rand der Länge nach über
das Auto. Der Wagen musste aus den 80ern sein. Ein altes, formschönes
Vehikel. Speziell, in einem autofreien Gebiet? Das Fahrzeug wirkte
verloren. Als hätte es jemand versehentlich hier abgestellt. Das Bild 
erinnerte Louis an einen Film, der ihm nicht mehr
einfallen wollte. Und sofort auch an die wechselnden, formschönen
Oldtimer in Vaters Autogarage.

Jetzt tauchte ein Mann im Türrahmen auf. Er kam die drei Treppenstufen
hinunter und schlich sich behutsam zur liegenden Frau auf der Bank. Er
beugte sich zu ihr hinunter und schien sie zu küssen. Die Frau streckte
die Arme aus und setzte sich auf. Das Lied startete von Neuem. Zum
vierten Mal. Der Mann zog die Frau zu sich und sie umarmten sich. Oder
tanzten sie?
Erstaunt legte Louis den Kopf schräg und kniff die Augen zusammen. Der
BMW in der Ecke des Bildes gab dieser kleinen Aufführung einen
speziellen Touch. Hinreissend skurril, dachte er. 

Mit dem Aufstieg auf die Bettmeralp öffnete sich ihm endgültig die Tür in eine fremde Welt. 
Er stand in einer Umgebung, die ihn sonderbar aufwühlte und
würde sich nun mit Entschlossenheit den Leuten vorstellen. Wenn sie ihn
nicht gebrauchen konnten, ginge seine Reise eben anders weiter. Louis
straffte sich, durchquerte die letzten Meter Wiese und zirkelte am Rande
eines Gemüsegartens direkt auf das Paar zu. Sie hatten sich in seine
Richtung umgedreht.
